% requer \usepackage{booktabs}. Tabela descritiva; sem interpretacao causal.
\begin{table}[htbp]\centering
\caption{Variacao acumulada desde 2010 dos componentes da idade de entrada e fracao atribuida a cada um}
\label{tab:T3b_cumulativa}
\small
\setlength{\tabcolsep}{4pt}
\begin{tabular}{lrrrrr}
\toprule
Ano & $\Delta$ Idade solic. & $\Delta$ Espera & $\Delta$ Idade entrada & Frac. idade solic. & Frac. espera \\
\midrule
2010 & 0.00 & 0.00 & 0.00 & n/d & n/d \\
2011 & -7.37 & 3.75 & -3.64 & 2.024 & -1.029 \\
2012 & -6.64 & 0.99 & -5.67 & 1.170 & -0.174 \\
2013 & -9.27 & 1.50 & -7.81 & 1.187 & -0.192 \\
2014 & -12.41 & 1.18 & -11.23 & 1.106 & -0.106 \\
2015 & -13.71 & -0.02 & -13.77 & 0.996 & 0.001 \\
2016 & -14.92 & -0.21 & -15.11 & 0.987 & 0.014 \\
2017 & -15.34 & -1.34 & -16.65 & 0.921 & 0.080 \\
2018 & -15.31 & -3.75 & -19.09 & 0.802 & 0.196 \\
\bottomrule
\end{tabular}
\\[3pt]{\footnotesize \textit{Note:} $\Delta$ = diferenca da media em relacao a 2010 (meses). Pela identidade, $\Delta$Idade entrada $=$ $\Delta$Idade solic. $+$ $\Delta$Espera; a fracao de cada componente e $\Delta$componente / $\Delta$Idade entrada. Decomposicao contabil, sem interpretacao causal.}
\end{table}
